% Phase-3 real switched-capacitor CMFB (replaces the ideal E1 VCVS).
% Continuous sense/hold caps C_C (=C7,C8) couple the output CM to the control
% node Vbcm (fast, ~diff GBW); switched caps C_CM (=Cs1,Cs2) refresh the DC level
% each clock cycle (sample on phi1 between Vcmr and Vbn, transfer on phi2 between
% Vout and Vbcm).  A weak ideal-DC helper (E through large R_help) sets the DC
% operating point for .op/.ac only.  Vbcm drives the output-CM sink gates Mn5,Mn7.
% Shared preamble for all CircuiTikZ schematics in this project.
% Each circuit file should be a self-contained `standalone` document that
% \input{}s this file.  Compile with:  pdflatex foo.tex
\documentclass[border=6pt]{standalone}
\usepackage{circuitikz}
\usepackage{amsmath}
\usetikzlibrary{calc,positioning}

% Project style — tuned for analog CMOS schematics
\ctikzset{
    bipoles/length=1.1cm,
    transistors/arrow pos=end,
    nodes width=0.06,
    logic ports=ieee,
}

% Default tikz options applied inside circuitikz environments
\tikzset{
    every picture/.style={
        line cap=round,
        line join=round,
        line width=0.6pt,
        font=\small,
    },
    every node/.style={font=\small},
}

% Convenience macros
\newcommand{\Vdd}{\ensuremath{V_{dd}}}
\newcommand{\Vddb}{\ensuremath{V_{ddb}}}
\newcommand{\Vss}{\ensuremath{V_{ss}}}

\begin{document}
\begin{circuitikz}[american, scale=1.0, line width=0.6pt]

  %% ===== control-node rail Vbcm =====
  \draw (12,3.2) -- (12,-3.2);
  \node[above] at (12,3.2) {$V_{bcm}$};
  \fill (12,0) circle (1.6pt);
  \draw[-latex] (12,0) -- ++(2.2,0) node[right,align=left]{to gates of\\$M_{n5},M_{n7}$\\(output CM sinks)};

  %% ===== output terminals =====
  \draw (0,3) node[left]{$V_{out+}$} to[short,o-] (1.2,3) coordinate (vop);
  \draw (0,-3) node[left]{$V_{out-}$} to[short,o-] (1.2,-3) coordinate (von);

  %% ===== continuous sense/hold caps C_C  (= C7, C8 = 1.25 pF) =====
  \draw (vop) to[C, l=$C_{C}$] (6.5,3) -- (12,3);
  \draw (von) to[C, l_=$C_{C}$] (6.5,-3) -- (12,-3);

  %% ===== + side switched sampling branch (C_CM = Cs1 = 0.5 pF) =====
  \draw (2.4,3) -- (2.4,1.4) coordinate (tp);
  \draw (tp) to[nos, l=$\phi_2$] (4.2,1.4) coordinate (ap);
  \draw (ap) to[C, l=$C_{CM}$] (8.2,1.4) coordinate (bp);
  \draw (bp) to[nos, l=$\phi_2$] (12,1.4);
  % phi1: ap -> Vcmr ,  bp -> Vbn
  \draw (ap) -- (4.2,2.3) to[nos, l=$\phi_1$] (6.0,2.3) node[right]{$V_{cmr}$};
  \draw (bp) -- (8.2,0.5) to[nos, l_=$\phi_1$] (10.0,0.5) node[right]{$V_{bn}$};

  %% ===== - side switched sampling branch (C_CM = Cs2 = 0.5 pF) =====
  \draw (2.4,-3) -- (2.4,-1.4) coordinate (tn);
  \draw (tn) to[nos, l_=$\phi_2$] (4.2,-1.4) coordinate (an);
  \draw (an) to[C, l_=$C_{CM}$] (8.2,-1.4) coordinate (bn);
  \draw (bn) to[nos, l_=$\phi_2$] (12,-1.4);
  \draw (an) -- (4.2,-2.3) to[nos, l_=$\phi_1$] (6.0,-2.3) node[right]{$V_{cmr}$};
  \draw (bn) -- (8.2,-0.5) to[nos, l=$\phi_1$] (10.0,-0.5) node[right]{$V_{bn}$};

  %% ===== weak ideal-DC CMFB helper (OP/AC only) =====
  % drawn as a labelled block: senses the output CM via large R1,R2 averaging,
  % compares to Vcmr, drives Vbcm through R_help (100 MOhm => kHz BW, DC only).
  \node[draw, rectangle, rounded corners, align=center, font=\footnotesize,
        minimum width=3.6cm, minimum height=1.3cm] (H) at (4.3,-5.0)
    {ideal DC CMFB ($E$)\\senses $V_{out}$ CM ($R_1,R_2$)\\ref.\ $V_{cmr}$};
  % CM sense input from the two outputs
  \draw (vop) -- (1.7,3) coordinate(t1) ;
  \draw (1.7,3) -- (1.7,-5.0) -- (H.west);
  \draw (von) -- (1.7,-3);
  \node[font=\footnotesize, anchor=east] at (1.55,-1.0) {$R_1,R_2$};
  % output through R_help up to the Vbcm rail
  \draw (H.east) to[R, l=$R_{help}$] (9.0,-5.0) -- (12,-5.0) -- (12,-3.2);
  \node[font=\footnotesize, anchor=west] at (9.2,-5.35) {$100\,$M$\Omega$};
  \node[align=left, font=\footnotesize] at (6.6,-6.2)
    {weak DC helper: sets the bias point for \texttt{.op}/\texttt{.ac} only\\(BW $\approx$ kHz $\ll$ signal band; the SC network does all in-band CM correction)};

  %% ===== clock note =====
  \node[align=left, font=\footnotesize] at (3.0,4.6)
    {$\phi_1$: sample / reset\\$\phi_2$: amplify (non-overlap, 20\,MHz)};

\end{circuitikz}
\end{document}
